\documentclass[11pt, letterpaper]{article}

\usepackage{amsmath}
\usepackage{amssymb}
\usepackage{amsthm}
\usepackage{geometry}

\geometry{
    top=1in,
    bottom=1in,
    left=1in,
    right=1in
}

\newcommand{\Prob}{\mathbb{P}}
\newcommand{\comp}[1]{\overline{#1}}

\title{\textbf{Homework 2 (Chapter 2 Part 1)}}
\author{
    Tashfeen Hatim Said Omran \\ 
    UCF ID: 5862894 \\ 
    Class: STA4321
}
\date{September 2026}

\begin{document}

\maketitle

\section*{Problem 2.5}
\textbf{a)} I use the given identities $A = A \cap S$ and $S = B \cup \comp{B}$. Substituting the latter into the former, I get $A = A \cap (B \cup \comp{B})$. Applying the distributive law yields:
\[ A = (A \cap B) \cup (A \cap \comp{B}) \]

\textbf{b)} If $B \subset A$, then all elements of $B$ are in $A$, meaning $A \cap B = B$. Substituting this into my result from part a gives:
\[ A = B \cup (A \cap \comp{B}) \]

\textbf{c)} To show they are mutually exclusive, I check their intersection:
\[ (A \cap B) \cap (A \cap \comp{B}) = A \cap (B \cap \comp{B}) = A \cap \emptyset = \emptyset \]
Since the intersection is empty, they are mutually exclusive. Thus, the identity from (a) expresses $A$ as the union of two mutually exclusive sets.

\textbf{d)} I check the intersection of $B$ and $A \cap \comp{B}$:
\[ B \cap (A \cap \comp{B}) = A \cap (B \cap \comp{B}) = A \cap \emptyset = \emptyset \]
This shows they are mutually exclusive. In part b, I found that if $B \subset A$, then $A = B \cup (A \cap \comp{B})$. Therefore, $A$ is the union of the two mutually exclusive sets $B$ and $(A \cap \comp{B})$.

\section*{Problem 2.8}
\textbf{a)} Definitions of my subsets:
\begin{itemize}
    \item $A = \{(x,y) \mid y \in \{2,4,6\}\}$
    \item $B = \{(x,y) \mid (x+y) \text{ is even}\}$
    \item $C = \{(x,y) \mid x \text{ is odd or } y \text{ is odd}\}$
\end{itemize}

\textbf{b)} I list the requested points:
\begin{itemize}
    \item \textbf{$A$}: $\{(1,2), (2,2), (3,2), (4,2), (5,2), (6,2), (1,4), (2,4), (3,4), (4,4), (5,4), (6,4), (1,6), (2,6), (3,6), (4,6), (5,6), (6,6)\}$
    \item \textbf{$\comp{C}$} (both even): $\{(2,2), (2,4), (2,6), (4,2), (4,4), (4,6), (6,2), (6,4), (6,6)\}$
    \item \textbf{$A \cap B$} (second even, sum even $\implies$ both even): Identical to $\comp{C}$.
    \item \textbf{$A \cap \comp{B}$} (second even, sum odd $\implies$ first odd): $\{(1,2), (3,2), (5,2), (1,4), (3,4), (5,4), (1,6), (3,6), (5,6)\}$
    \item \textbf{$\comp{A} \cup B$} (second odd or sum even): This contains 27 points. It includes all points where both are odd (9 points), both are even (9 points), and first even/second odd (9 points).
    \item \textbf{$\comp{A} \cap C$} (second odd and at least one odd): Since the second die being odd guarantees at least one is odd, this is just $\comp{A}$.
\end{itemize}

\section*{Problem 2.11}
\textbf{a)} I know the sum of all simple event probabilities is $1$.
\[ \Prob(E_1) + \Prob(E_2) + \Prob(E_3) + \Prob(E_4) + \Prob(E_5) = 1 \]
\[ 0.15 + 0.15 + 0.40 + 2\Prob(E_5) + \Prob(E_5) = 1 \]
\[ 0.70 + 3\Prob(E_5) = 1 \implies 3\Prob(E_5) = 0.30 \implies \Prob(E_5) = 0.10 \]
Since $\Prob(E_4) = 2\Prob(E_5)$, I get $\Prob(E_4) = 0.20$.

\textbf{b)} I am given $\Prob(E_1) = 3\Prob(E_2) = 0.3$. So, $\Prob(E_1) = 0.3$ and $\Prob(E_2) = 0.1$. The remaining three events have equal probability $x$.
\[ 0.3 + 0.1 + 3x = 1 \implies 0.4 + 3x = 1 \implies 3x = 0.6 \implies x = 0.2 \]
The remaining probabilities are $\Prob(E_3) = \Prob(E_4) = \Prob(E_5) = 0.2$.

\section*{Problem 2.19}
\textbf{a)} I denote the choices over the two days as ordered pairs. The sample space is:
\[ S = \{(V_1,V_1), (V_1,V_2), (V_1,V_3), (V_2,V_1), (V_2,V_2), (V_2,V_3), (V_3,V_1), (V_3,V_2), (V_3,V_3)\} \]
\textbf{b)} Because selection is random, each of the 9 points is equally likely. I assign a probability of $1/9$ to each sample point. \\
\textbf{c)} I define my events and compute the probabilities:
\begin{itemize}
    \item $A = \{(V_1,V_1), (V_2,V_2), (V_3,V_3)\} \implies \Prob(A) = 3/9 = 1/3$
    \item $B = \{(V_1,V_2), (V_2,V_1), (V_2,V_2), (V_2,V_3), (V_3,V_2)\} \implies \Prob(B) = 5/9$
    \item $A \cap B = \{(V_2,V_2)\} \implies \Prob(A \cap B) = 1/9$
    \item $\Prob(A \cup B) = \Prob(A) + \Prob(B) - \Prob(A \cap B) = 3/9 + 5/9 - 1/9 = 7/9$
\end{itemize}

\section*{Problem 2.29}
\textbf{a)} The experiment is randomly selecting 2 people from a pool of 6 (2 women, 4 men). The order does not matter. One sample point would be selecting woman 1 and man 2, denoted as $\{W_1, M_2\}$. \\
\textbf{b)} I list out all $\binom{6}{2} = 15$ possible combinations for the sample space:
\[ S = \left\{ \begin{array}{l} \{W_1,W_2\}, \{W_1,M_1\}, \{W_1,M_2\}, \{W_1,M_3\}, \{W_1,M_4\}, \\ \{W_2,M_1\}, \{W_2,M_2\}, \{W_2,M_3\}, \{W_2,M_4\}, \\ \{M_1,M_2\}, \{M_1,M_3\}, \{M_1,M_4\}, \{M_2,M_3\}, \{M_2,M_4\}, \{M_3,M_4\} \end{array} \right\} \]
\textbf{c)} The event where both jurors are women is exactly 1 sample point: $\{W_1, W_2\}$. The probability is $1/15$.

\section*{Problem 2.38}
I use the basic multiplication principle. I multiply the number of choices for each course:
\[ 4 \text{ (appetizers)} \times 3 \text{ (salads)} \times 4 \text{ (entrees)} \times 5 \text{ (desserts)} = 240 \text{ different dinners} \]

\section*{Problem 2.42}
Because the positions are "distinctly different," order matters. I need to calculate permutations for choosing 3 out of 10 engineers:
\[ P(10, 3) = \frac{10!}{(10-3)!} = 10 \times 9 \times 8 = 720 \text{ ways} \]

\section*{Problem 2.43}
This is a partitioning problem where I distribute 9 items into groups of size 3, 5, and 1. I apply the multinomial coefficient:
\[ \binom{9}{3, 5, 1} = \frac{9!}{3! \cdot 5! \cdot 1!} = \frac{9 \times 8 \times 7 \times 6}{3 \times 2 \times 1} = 504 \text{ distinct ways} \]

\section*{Problem 2.44}
\textbf{a)} There are a total of 9 dispatch slots. Airport C receives exactly 1 slot. The probability that the single broken taxi is assigned to C is simply:
\[ \Prob(\text{to C}) = \frac{1}{9} \]
\textbf{b)} There are $\binom{9}{3} = 84$ total ways to place the 3 indistinguishable broken taxis into the 9 slots. For every airport to receive exactly one, I must pick 1 slot from A's 3 slots, 1 from B's 5 slots, and 1 from C's 1 slot. This gives $\binom{3}{1} \times \binom{5}{1} \times \binom{1}{1} = 15$ favorable ways.
\[ \Prob(\text{one to each}) = \frac{15}{84} = \frac{5}{28} \]

\section*{Problem 2.51}
This is a hypergeometric distribution problem. I am drawing 3 winners from a population of 50 tickets. The 4 organizers hold the "success" tickets. The total possible combinations of winning draws is $\binom{50}{3} = 19600$.

\textbf{a) All of the prizes:}
\[ \Prob(X = 3) = \frac{\binom{4}{3}\binom{46}{0}}{19600} = \frac{4 \times 1}{19600} = \frac{1}{4900} \]

\textbf{b) Exactly two of the prizes:}
\[ \Prob(X = 2) = \frac{\binom{4}{2}\binom{46}{1}}{19600} = \frac{6 \times 46}{19600} = \frac{276}{19600} = \frac{69}{4900} \]

\textbf{c) Exactly one of the prizes:}
\[ \Prob(X = 1) = \frac{\binom{4}{1}\binom{46}{2}}{19600} = \frac{4 \times 1035}{19600} = \frac{4140}{19600} = \frac{207}{980} \]

\textbf{d) None of the prizes:}
\[ \Prob(X = 0) = \frac{\binom{4}{0}\binom{46}{3}}{19600} = \frac{1 \times 15180}{19600} = \frac{15180}{19600} = \frac{759}{980} \]

\end{document}